% !TeX program = lualatex

% ======================================================================
% PRODUCCIÓN DE DATOS CUALITATIVOS II: ENTREVISTA E INFORMACIÓN SECUNDARIA
% (BEAMER, THEME METROPOLIS)
% ----------------------------------------------------------------------
% Motor de compilación requerido: LuaLaTeX
%   lualatex 16-produccionDatosCual2.tex
%   biber    16-produccionDatosCual2
%   lualatex 16-produccionDatosCual2.tex  (x2, para resolver referencias)
% NO compilar con pdflatex: fontspec y polyglossia no son compatibles.
%
% Semana 11 (Unidad 3), segunda sesión. Primera parte: la entrevista
% cualitativa (Piovani, en Marradi et al. 2007, cap. 12; Bryman et al.
% 2021, cap. 19; Vieytes 2004, cap. 23): definición, cuándo entrevistar
% y no observar, procedimientos (selección de Gorden, guía en embudo,
% nueve tipos de preguntas de Kvale, tácticas, registro), sesgos y cómo
% reducirlos, variantes (Fideli y Marradi; focalizada, historia de vida,
% historia oral, grupal, élites, en línea). Segunda parte: información
% secundaria cualitativa (Bryman et al. 2021, cap. 22; Scott 1990;
% Vieytes 2004, cap. 6; Heaton, en Alasuutari et al. 2008): tipos de
% documentos, procedimiento, criterios de Scott, recomendaciones cuando
% solo hay datos secundarios o cuando complementan, y procedimientos
% para evitar sesgos (supervivencia selectiva, silencios del archivo,
% autoría y destinatario). Incluye ejemplos Q1 y dos ejercicios breves.
% Fuentes: notebook NotebookLM "LogicaMetodos" y artículos Q1.
%
% Duración estimada: ~90 min
%   Dónde estamos e índice                                        ~2 min
%   1. La entrevista: qué es y cuándo usarla                       ~8 min
%   2. Procedimientos y técnicas de entrevista                    ~14 min
%   3. Sesgos en la entrevista                                     ~7 min
%   4. Variantes de entrevista                                     ~7 min
%   5. Ejercicio: preguntas de entrevista
%   6. Información secundaria cualitativa: procedimientos         ~12 min
%   7. Solo secundaria o como complemento                          ~7 min
%   8. Sesgos en datos secundarios                                 ~7 min
%   9. Ejercicio: evaluar un documento
%   Cierre                                                         ~3 min
% ======================================================================

\documentclass[12pt,aspectratio=169]{beamer}

% ---------------------------------------------------------------------
% 1) TEMA: METROPOLIS
% ---------------------------------------------------------------------
\usetheme{metropolis}

% ---------------------------------------------------------------------
% 2) IDIOMA (requiere LuaLaTeX o XeLaTeX)
% ---------------------------------------------------------------------
\usepackage{polyglossia}
\setmainlanguage{spanish}

% ---------------------------------------------------------------------
% 3) COLORES INSTITUCIONALES (Universidad Torcuato Di Tella)
% ---------------------------------------------------------------------
\usepackage[dvipsnames]{xcolor}
\definecolor{utdtYellow}{HTML}{F9BC00}
\definecolor{utdtDark}{HTML}{1B1C20}
\definecolor{utdtGray}{HTML}{808083}

\setbeamercolor{normal text}{fg=utdtDark, bg=white}
\setbeamercolor{alerted text}{fg=utdtYellow}
\setbeamercolor{example text}{fg=utdtGray}

% ---------------------------------------------------------------------
% 4) OPCIONES DEL TEMA
% ---------------------------------------------------------------------
\metroset{
  progressbar=frametitle,
  sectionpage=progressbar,
  numbering=fraction,
  block=fill
}

\usepackage{booktabs}
\usepackage{graphicx}
\usepackage{pifont}
\usepackage{amsmath}

% ---------------------------------------------------------------------
% 5) TIKZ (diagramas de la clase)
% ---------------------------------------------------------------------
\usepackage{tikz}
\usetikzlibrary{arrows.meta, positioning, shapes.geometric, fit, backgrounds, calc, decorations.pathreplacing}
\usepackage{adjustbox}

% Estilos comunes de esquemas conceptuales (cajas)
\tikzset{
  arr/.style={-{Stealth[length=2.8mm]}, thick, draw=utdtDark},
  badarr/.style={-{Stealth[length=2.8mm]}, thick, draw=BrickRed, dashed},
  bx/.style={draw=utdtDark, thick, rounded corners, align=center, fill=white, font=\scriptsize, inner sep=4pt},
  bxk/.style={bx, fill=utdtYellow},
  bxg/.style={draw=utdtGray, thick, rounded corners, align=center, fill=utdtGray!10, font=\tiny, inner sep=3pt},
  bxr/.style={bx, draw=BrickRed},
  lbl/.style={font=\tiny, text=utdtDark, fill=white, inner sep=1.5pt, align=center},
  glbl/.style={font=\tiny, text=utdtGray, inner sep=1pt, align=center},
  garr/.style={-{Stealth[length=2.5mm]}, thick, draw=utdtGray},
  darr/.style={-{Stealth[length=2.5mm]}, thick, dashed, draw=utdtGray},
  biarr/.style={{Stealth[length=2.4mm]}-{Stealth[length=2.4mm]}, thick, draw=utdtDark},
  doc/.style={draw=utdtDark, thick, fill=white, font=\tiny, align=left, inner sep=5pt},
}

% ---- Datos editables --------------------------------------------------
\newcommand{\NombreProfesor}{Valentín Vergara Hidd}
\newcommand{\NombreCurso}{Lógica y Técnicas de Investigación Social}
\newcommand{\Departamento}{Estudios Históricos y Sociales}
\newcommand{\CorreoProfesor}{valentin.vergara@utdt.edu}
\newcommand{\ArchivoLogo}{logoColor.png}

\newcommand{\DeterminarSemestre}[1]{%
  \ifnum#1<7
    Primer Semestre%
  \else
    Segundo Semestre%
  \fi
}
\newcommand{\Semestre}{\DeterminarSemestre{\month} \the\year}
% -----------------------------------------------------------------------
%% REFERENCIAS
\usepackage[style=apa, backend=biber]{biblatex}
\addbibresource{referencias.bib}

\title{\NombreCurso}
\subtitle{Datos cualitativos II: entrevista e información secundaria}
\author{\NombreProfesor}
\institute{\Departamento \\ Universidad Torcuato Di Tella}
\date{\Semestre. {\small \today}}
\titlegraphic{\includegraphics[height=0.8cm]{\ArchivoLogo}\vspace*{0.5cm}}

\hypersetup{
  pdftitle={\NombreCurso -- Producción de datos cualitativos II: entrevista e información secundaria},
  pdfauthor={\NombreProfesor}
}

\begin{document}

\maketitle

% ============================ APERTURA ==============================
\begin{frame}[t]{Dónde estamos: del hacer al decir y a las huellas}
  \begin{center}
  \adjustbox{max width=0.97\textwidth, max height=0.6\textheight}{%
  \begin{tikzpicture}
    \node[bx, fill=utdtDark, text=white, font=\scriptsize\bfseries, text width=2.4cm, minimum height=1.0cm] (p) at (0,0) {pregunta de investigación};
    \node[bxg, font=\scriptsize, text width=3.2cm, minimum height=1.3cm] (o) at (0,2.3) {\textbf{observación} (clase anterior)\\lo que la gente \textbf{hace}};
    \node[bxk, text width=3.2cm, minimum height=1.3cm] (e) at (-4.8,-1.6) {\textbf{entrevista} (hoy, parte 1)\\lo que la gente \textbf{dice}};
    \node[bxk, text width=3.2cm, minimum height=1.3cm] (d) at (4.8,-1.6) {\textbf{información secundaria} (hoy, parte 2)\\las \textbf{huellas} que otros dejaron};
    \draw[garr] (p) -- (o); \draw[arr] (p) -- (e); \draw[arr] (p) -- (d);
    \draw[biarr, dashed, draw=utdtGray] (e.east) -- node[lbl]{triangulación} (d.west);
    \node[glbl, font=\scriptsize] at (0,-2.9) {semana 12: \textbf{analizar} lo producido (codificación y teoría fundamentada)};
  \end{tikzpicture}}
  \end{center}
\end{frame}

\begin{frame}{Temas de hoy}
  \footnotesize
  \tableofcontents
\end{frame}

% ============================ SECCIÓN 1 ============================
\section{La entrevista: qué es y cuándo usarla}

\begin{frame}[t]{Una conversación con propósito}
  \begin{center}
  \adjustbox{max width=0.97\textwidth, max height=0.6\textheight}{%
  \begin{tikzpicture}
    \node[bxk, text width=4.4cm, minimum height=1.8cm] (def) at (0,0) {\textbf{entrevista en profundidad}\\conversación entre dos personas, \textbf{dirigida y registrada} por el investigador, para producir un \textbf{discurso continuo} sobre un tema de la investigación};
    \node[bx, text width=2.8cm] (v1) at (-5.6,1.6) {no un registro de hechos, sino un \textbf{``decir sobre el hacer''}};
    \node[bx, text width=2.8cm] (v2) at (-5.6,-1.6) {el entrevistado es el \textbf{experto}; el entrevistador ayuda a verbalizar (rol ``mayéutico'')};
    \node[bx, text width=2.8cm] (c1) at (5.6,1.6) {\textbf{saberes explícitos}: quién pregunta, para qué; consentimiento informado};
    \node[bx, text width=2.8cm] (c2) at (5.6,-1.6) {\textbf{saberes implícitos}: códigos lingüísticos y culturales compartidos};
    \draw[arr] (v1) -- (def); \draw[arr] (v2) -- (def);
    \draw[arr] (c1) -- (def); \draw[arr] (c2) -- (def);
    \node[font=\scriptsize\bfseries] at (5.6,2.7) {contrato comunicativo};
    \node[font=\scriptsize\bfseries] at (-5.6,2.7) {qué produce};
  \end{tikzpicture}}
  \end{center}
  \begin{center}
  \tiny Alonso (1998) y Montesperelli (1998), en Piovani, \textcite[cap.~12, pp.~216--220]{marradi2007metodologia}. Cerca del 90\,\% de los estudios empíricos usan algún tipo de entrevista (Fideli y Marradi, 1996).
  \end{center}
\end{frame}

\begin{frame}[t]{Una estructura paradójica}
  \begin{center}
  \adjustbox{max width=0.9\textwidth, max height=0.56\textheight}{%
  \begin{tikzpicture}
    \foreach \y/\a/\b in {1.6/intimidad/impersonalidad, 0.4/sociabilidad/profesionalismo, -0.8/conversación natural/agenda del investigador, -2.0/empatía/no juzgar ni opinar} {
      \node[bx, text width=2.8cm, fill=utdtYellow!50] (a) at (-3.6,\y) {\a};
      \node[bx, text width=2.8cm] (b) at (3.6,\y) {\b};
      \draw[biarr, draw=utdtGray] (a) -- (b);
    }
    \node[bxg, font=\scriptsize, text width=3.6cm] at (0,3.0) {intercambio breve, instrumental, entre extraños};
  \end{tikzpicture}}
  \end{center}
  \begin{center}
  \tiny Rosenblum (1987), en \textcite[p.~218]{marradi2007metodologia}. El entrevistador sostiene ambas cosas a la vez: por eso la entrevista se describe como un \textbf{oficio} más que como una receta.
  \end{center}
\end{frame}

\begin{frame}[t]{¿Cuándo entrevistar y no observar?}
  \begin{center}
  \adjustbox{max width=0.98\textwidth, max height=0.66\textheight}{%
  \begin{tikzpicture}[c/.style={bx, text width=3.0cm, minimum height=1.2cm}]
    \node[bxk, circle, minimum size=2.0cm, text width=1.8cm, font=\scriptsize\bfseries] (o) at (0,0) {entrevistar};
    \node[c] (c1) at (-5.0,1.9) {temas \textbf{no observables}: motivos, emociones, representaciones};
    \node[c] (c2) at (0,2.6) {\textbf{reconstruir} hechos pasados y trayectorias};
    \node[c] (c3) at (5.0,1.9) {temas \textbf{sensibles} donde observar sería invasivo o poco ético};
    \node[c] (c4) at (-5.0,-1.9) {mayor \textbf{cobertura}: personas y contextos dispersos};
    \node[c] (c5) at (0,-2.6) {foco \textbf{específico}; diseños \textbf{longitudinales} con reentrevistas};
    \node[c] (c6) at (5.0,-1.9) {comparar cómo distintas personas \textbf{dan sentido} a lo mismo};
    \foreach \c in {c1,c2,c3,c4,c5,c6} { \draw[arr] (\c) -- (o); }
  \end{tikzpicture}}
  \end{center}
  \begin{center}
  \tiny \textcite[sec.~19.7]{brymanetal2021socialresearch}; Alonso (1998) y Valles (1997), en \textcite[p.~220]{marradi2007metodologia}; \textcite{lamontswidler2014pluralism}.
  \end{center}
\end{frame}

\begin{frame}[t]{Lo que la entrevista no entrega}
  \begin{center}
  \adjustbox{max width=0.97\textwidth, max height=0.58\textheight}{%
  \begin{tikzpicture}[r/.style={bxr, text width=3.1cm, minimum height=1.2cm}]
    \node[bx, fill=utdtDark, text=white, font=\scriptsize\bfseries, text width=2.4cm] (e) at (0,0) {relato verbal en una situación artificial};
    \node[r] (r1) at (-5.0,1.4) {\textbf{validez ecológica}: no se ve el fenómeno en su ambiente (Cicourel)};
    \node[r] (r2) at (5.0,1.4) {lo \textbf{dado por sentado} no aparece en el relato};
    \node[r] (r3) at (-5.0,-1.4) {lo que se \textbf{dice} puede no ser lo que se \textbf{hace}: el ejecutivo ``participativo'' que decide solo};
    \node[r] (r4) at (5.0,-1.4) {preguntar ``\textbf{por qué}'' invita a justificaciones construidas después};
    \foreach \r in {r1,r2,r3,r4} { \draw[badarr] (e) -- (\r); }
  \end{tikzpicture}}
  \end{center}
  \begin{center}
  \tiny Piovani, en \textcite[pp.~220--221]{marradi2007metodologia}; \textcite[cap.~23]{vieytes2004metodologia}; \textcite{jerolmackkhan2014talk}. \textcite{smallcook2023why} muestran cómo entrevistar sobre motivos sin caer en racionalizaciones: pedir relatos concretos y secuencias, no explicaciones abstractas.
  \end{center}
\end{frame}

% ============================ SECCIÓN 2 ============================
\section{Procedimientos y técnicas para entrevistar}

\begin{frame}[t]{La entrevista, en secuencia}
  \begin{center}
  \adjustbox{max width=0.98\textwidth, max height=0.6\textheight}{%
  \begin{tikzpicture}[s/.style={bx, text width=1.9cm, minimum height=1.7cm, font=\tiny}]
    \node[s, fill=utdtYellow] (a) at (0,0) {\textbf{1. seleccionar}\\a quiénes y cuántos};
    \node[s] (b) at (2.35,0) {\textbf{2. guía}\\temas en embudo; entrevistas piloto};
    \node[s] (c) at (4.7,0) {\textbf{3. contacto}\\propósito, duración, confidencialidad};
    \node[s] (d) at (7.05,0) {\textbf{4. preparar}\\lugar, registro, presentación personal};
    \node[s] (e) at (9.4,0) {\textbf{5. conducir}\\preguntas y tácticas};
    \node[s] (f) at (11.75,0) {\textbf{6. cerrar}\\agradecer; post-entrevista};
    \node[s] (g) at (14.1,0) {\textbf{7. transcribir}\\y anotar el contexto};
    \foreach \x/\y in {a/b, b/c, c/d, d/e, e/f, f/g} { \draw[arr] (\x) -- (\y); }
    \draw[darr] (g.south) to[out=-150, in=-30] node[lbl, below]{lo aprendido en cada entrevista refina la guía y la selección} (b.south);
  \end{tikzpicture}}
  \end{center}
  \begin{center}
  \tiny \textcite[pp.~221--225]{marradi2007metodologia}; \textcite[secs.~19.3--19.4]{brymanetal2021socialresearch}; plan de entrevistas semiestructuradas de \textcite[cap.~14]{vieytes2004metodologia}.
  \end{center}
\end{frame}

\begin{frame}[t]{Seleccionar a quiénes entrevistar}
  \begin{columns}[T]
    \begin{column}{0.5\textwidth}
      \begin{center}
      \adjustbox{max width=\textwidth, max height=0.62\textheight}{%
      \begin{tikzpicture}
        \fill[utdtYellow!20, draw=utdtDark, thick] (-3.6,0) -- (3.6,0) -- (3.0,-1.0) -- (-3.0,-1.0) -- cycle;
        \fill[utdtYellow!45, draw=utdtDark, thick] (-3.0,-1.0) -- (3.0,-1.0) -- (2.4,-2.0) -- (-2.4,-2.0) -- cycle;
        \fill[utdtYellow!70, draw=utdtDark, thick] (-2.4,-2.0) -- (2.4,-2.0) -- (1.8,-3.0) -- (-1.8,-3.0) -- cycle;
        \fill[utdtYellow, draw=utdtDark, thick] (-1.8,-3.0) -- (1.8,-3.0) -- (1.2,-4.0) -- (-1.2,-4.0) -- cycle;
        \node[font=\tiny, align=center] at (0,-0.5) {1. ¿quiénes \textbf{tienen la información}?};
        \node[font=\tiny, align=center] at (0,-1.5) {2. ¿quiénes son \textbf{accesibles}?};
        \node[font=\tiny, align=center] at (0,-2.5) {3. ¿quiénes están \textbf{dispuestos}?};
        \node[font=\tiny, align=center, text width=2.2cm] at (0,-3.5) {4. ¿quiénes \textbf{comunican} bien?};
      \end{tikzpicture}}
      \end{center}
      \tiny Gorden (1969), según Valles, en \textcite[p.~222]{marradi2007metodologia}.
    \end{column}
    \begin{column}{0.48\textwidth}
      \begin{center}
      \adjustbox{max width=\textwidth, max height=0.64\textheight}{%
      \begin{tikzpicture}[s/.style={bx, text width=4.4cm, minimum height=0.65cm, font=\tiny}]
        \node[s, fill=utdtDark, text=white] (t) at (0,0) {\textbf{estrategias de selección}};
        \node[s] (a) at (0,-0.85) {\textbf{intencional}: segmentar por criterios relevantes y cubrir la diversidad de experiencias};
        \node[s] (b) at (0,-1.85) {\textbf{bola de nieve}: poblaciones poco visibles o dispersas};
        \node[s] (c) at (0,-2.75) {\textbf{oportunista}: acceso garantizado (con riesgos)};
        \node[s, fill=utdtYellow] (d) at (0,-3.75) {¿cuántas? \textbf{saturación}: hasta que nuevas entrevistas no aportan elementos nuevos};
        \foreach \x/\y in {t/a, a/b, b/c} { \draw[garr] (\x) -- (\y); }
        \draw[arr] (c) -- (d);
      \end{tikzpicture}}
      \end{center}
      \tiny \textcite[pp.~222--223]{marradi2007metodologia}; \textcite{henninkkaiser2022saturation}; \textcite{malterud2016informationpower}.
    \end{column}
  \end{columns}
\end{frame}

\begin{frame}[t]{La guía: un recordatorio, no un interrogatorio}
  \begin{center}
  \adjustbox{max width=0.97\textwidth, max height=0.6\textheight}{%
  \begin{tikzpicture}
    \node[bx, text width=2.6cm, minimum height=1.1cm] (i) at (0,0) {\textbf{inicio}\\preguntas amplias y abiertas};
    \node[bxk, text width=2.6cm, minimum height=1.1cm] (m) at (4.2,0) {\textbf{intermedias}\\experiencias concretas, sentimientos};
    \node[bx, text width=2.6cm, minimum height=1.1cm] (f) at (8.4,0) {\textbf{cierre}\\balance, cambios, consejos; ¿algo más?};
    \draw[arr] (i) -- (m); \draw[arr] (m) -- (f);
    \draw[decorate, decoration={brace, amplitude=6pt}, thick] (-1.3,0.9) -- (9.7,0.9) node[midway, above=8pt, font=\scriptsize, text width=10cm, align=center]{\textbf{embudo}: de lo general a lo particular, de lo impersonal a lo personal, de lo informativo a lo interpretativo};
    \node[bxg, font=\scriptsize, text width=3.0cm] (g1) at (0.6,-2.0) {a mayor conocimiento del tema, guía más detallada};
    \node[bxg, font=\scriptsize, text width=3.0cm] (g2) at (4.2,-2.0) {preguntas \textbf{directas} hacia el final, para no orientar el resto};
    \node[bxg, font=\scriptsize, text width=3.0cm] (g3) at (7.8,-2.0) {se ajusta con entrevistas \textbf{exploratorias}};
  \end{tikzpicture}}
  \end{center}
  \begin{center}
  \tiny Tipos de preguntas por etapa de Charmaz (2002), en \textcite[sec.~19.3]{brymanetal2021socialresearch}. ``La entrevista debe seguir su propia dinámica'' (Combessie), en \textcite[p.~223]{marradi2007metodologia}; técnica del embudo en \textcite[cap.~14]{vieytes2004metodologia}.
  \end{center}
\end{frame}

\begin{frame}[t]{Nueve tipos de preguntas}
  \begin{center}
  \adjustbox{max width=0.98\textwidth, max height=0.7\textheight}{%
  \begin{tikzpicture}[k/.style={bxk, font=\tiny\bfseries, text width=2.0cm, minimum height=0.95cm}, e/.style={bx, font=\tiny, text width=2.4cm, minimum height=0.95cm, align=left}]
    \foreach \i/\t/\x in {
      0/{1. introductorias}/{``Cuénteme cómo empezó a trabajar mientras estudiaba.''},
      1/{2. de seguimiento}/{``¿Sí\ldots?''; repetir una palabra clave del relato},
      2/{3. de sondeo}/{``Dijo que prefiere no pedir ayuda, ¿qué lo hizo desistir?''}} {
      \node[k] at (0,{-1.15*\i}) {\t}; \node[e] at (2.4,{-1.15*\i}) {\x};
    }
    \foreach \i/\t/\x in {
      0/{4. de especificación}/{``¿Qué hizo entonces? ¿Cómo reaccionó su jefe?''},
      1/{5. directas}/{``¿Ha pensado en dejar la carrera?'' (mejor al final)},
      2/{6. indirectas}/{``¿Qué opinan sus compañeros de\ldots? ¿Y usted?''}} {
      \node[k] at (5.4,{-1.15*\i}) {\t}; \node[e] at (7.8,{-1.15*\i}) {\x};
    }
    \foreach \i/\t/\x in {
      0/{7. estructurantes}/{``Ahora quisiera pasar a otro tema, ¿le parece?''},
      1/{8. silencio}/{pausas que dan tiempo para pensar y ampliar},
      2/{9. interpretativas}/{``¿Quiere decir que el trabajo pasó a ser lo primero?''}} {
      \node[k] at (10.8,{-1.15*\i}) {\t}; \node[e] at (13.2,{-1.15*\i}) {\x};
    }
  \end{tikzpicture}}
  \end{center}
  \begin{center}
  \tiny \textcite{kvale1996interviews}, en \textcite[sec.~19.3]{brymanetal2021socialresearch}. La mayoría de las entrevistas semiestructuradas combinan casi todos estos tipos.
  \end{center}
\end{frame}

\begin{frame}[t]{Tácticas del entrevistador y cualidades que las sostienen}
  \begin{columns}[T]
    \begin{column}{0.48\textwidth}
      \begin{center}
      \adjustbox{max width=\textwidth, max height=0.62\textheight}{%
      \begin{tikzpicture}
        \node[bxk, circle, minimum size=1.6cm, text width=1.4cm, font=\tiny\bfseries] (t) at (0,0) {tácticas};
        \foreach \a/\x in {90/silencio, 40/{animación y elaboración}, -10/{relanzamiento: repetir una palabra}, -60/recapitulación, -120/aclaración, -170/{cambio de tema}, 140/{post-entrevista: ``micrófono cerrado''}} {
          \node[bx, font=\tiny, text width=1.9cm] (n\a) at (\a:2.5) {\x};
          \draw[garr] (t) -- (n\a);
        }
      \end{tikzpicture}}
      \end{center}
      \tiny Valles (1997) y Combessie (1996), en \textcite[p.~225]{marradi2007metodologia}.
    \end{column}
    \begin{column}{0.5\textwidth}
      \begin{center}
      \adjustbox{max width=\textwidth, max height=0.62\textheight}{%
      \begin{tikzpicture}[c/.style={bx, font=\tiny, text width=1.75cm, minimum height=0.75cm}]
        \foreach \i/\t in {0/{conocedor}, 1/{estructurante}, 2/{claro}, 3/{amable}} { \node[c] at (0,{-0.9*\i}) {\t}; }
        \foreach \i/\t in {0/{sensible}, 1/{abierto}, 2/{orientador}, 3/{crítico}} { \node[c] at (2.1,{-0.9*\i}) {\t}; }
        \foreach \i/\t in {0/{memorioso}, 1/{intérprete}, 2/{equilibrado}, 3/{éticamente sensible}} { \node[c, fill=utdtYellow!40] at (4.2,{-0.9*\i}) {\t}; }
        \node[bxg, font=\tiny, text width=5.6cm] at (2.1,-3.8) {base común: \textbf{escuchar}, ser flexible y \textbf{no juzgar}; no mostrar acuerdo ni desacuerdo};
      \end{tikzpicture}}
      \end{center}
      \tiny Diez cualidades de \textcite{kvale1996interviews} y dos agregadas por \textcite[sec.~19.4]{brymanetal2021socialresearch}.
    \end{column}
  \end{columns}
\end{frame}

\begin{frame}[t]{Lugar, registro y transcripción}
  \begin{center}
  \adjustbox{max width=0.97\textwidth, max height=0.6\textheight}{%
  \begin{tikzpicture}
    \node[bxk, text width=3.0cm, minimum height=1.4cm] (l) at (0,0) {\textbf{lugar}\\condiciona lo que se puede decir};
    \node[bx, text width=3.0cm, minimum height=1.4cm] (g) at (5.0,0) {\textbf{grabación}\\con consentimiento; libera la atención para escuchar y observar gestos};
    \node[bx, text width=3.0cm, minimum height=1.4cm] (t) at (10.0,0) {\textbf{transcripción}\\cada hora de audio exige varias horas de trabajo};
    \draw[arr] (l) -- (g); \draw[arr] (g) -- (t);
    \node[bxg, font=\tiny, text width=3.0cm] (l2) at (0,-2.0) {un adolescente no dice lo mismo en el patio del colegio, en una fiesta o en su casa frente a sus padres};
    \node[bxg, font=\tiny, text width=3.0cm] (g2) at (5.0,-2.0) {tomar notas a mano también es reactivo y corta la fluidez};
    \node[bxg, font=\tiny, text width=3.0cm] (t2) at (10.0,-2.0) {sumar notas de contexto: clima, interrupciones, lenguaje corporal, post-entrevista};
    \draw[garr] (l2) -- (l); \draw[garr] (g2) -- (g); \draw[garr] (t2) -- (t);
  \end{tikzpicture}}
  \end{center}
  \begin{center}
  \tiny \textcite[pp.~224--225]{marradi2007metodologia}; \textcite[sec.~19.4]{brymanetal2021socialresearch}. Lugar tranquilo y privado, para que el entrevistado no tema ser escuchado.
  \end{center}
\end{frame}

% ============================ SECCIÓN 3 ============================
\section{¿Cómo mantener libre de sesgos una entrevista?}

\begin{frame}[t]{Fuentes de sesgo en la entrevista}
  \begin{center}
  \adjustbox{max width=0.98\textwidth, max height=0.64\textheight}{%
  \begin{tikzpicture}[r/.style={bxr, text width=3.2cm, minimum height=1.2cm}]
    \node[bx, fill=utdtDark, text=white, font=\scriptsize\bfseries, text width=2.4cm] (o) at (0,0) {la interacción de entrevista};
    \node[r] (r1) at (-5.0,1.7) {\textbf{preguntas dirigidas}: la formulación sugiere la respuesta};
    \node[r] (r2) at (0,2.4) {\textbf{deseabilidad social}: decir lo aceptable, no lo que se piensa o hace};
    \node[r] (r3) at (5.0,1.7) {\textbf{efecto entrevistador}: edad, género, estatus y origen modelan el relato};
    \node[r] (r4) at (-5.0,-1.7) {\textbf{selección}: solo hablan los dispuestos y accesibles};
    \node[r] (r5) at (0,-2.4) {\textbf{memoria}: olvidos y reconstrucción del pasado desde el presente};
    \node[r] (r6) at (5.0,-1.7) {\textbf{interpretación impuesta}: el investigador pone palabras en boca del otro};
    \foreach \r in {r1,r2,r3,r4,r5,r6} { \draw[badarr] (o) -- (\r); }
  \end{tikzpicture}}
  \end{center}
  \begin{center}
  \tiny \textcite[cap.~23]{vieytes2004metodologia}; \textcite[secs.~19.4 y 19.6]{brymanetal2021socialresearch}; \textcite[pp.~220--224]{marradi2007metodologia}.
  \end{center}
\end{frame}

\begin{frame}[t]{Riesgos y procedimientos}
  \begin{center}
  \adjustbox{max width=0.97\textwidth, max height=0.7\textheight}{%
  \begin{tikzpicture}[r/.style={bxr, text width=3.0cm, minimum height=0.8cm, font=\tiny}, s/.style={bxk, text width=7.0cm, minimum height=0.8cm, font=\tiny}]
    \node[font=\scriptsize\bfseries] at (0,0.75) {riesgo};
    \node[font=\scriptsize\bfseries] at (6.2,0.75) {procedimiento};
    \foreach \i/\a/\b in {
      0/{preguntas dirigidas}/{preguntas abiertas y neutras; sin adjetivos ni supuestos; probar la guía en entrevistas piloto},
      1/{deseabilidad social}/{garantizar confidencialidad; preguntas indirectas o en tercera persona; pedir episodios concretos},
      2/{efecto entrevistador}/{ejercicio autorreflexivo previo; adaptar presentación y lenguaje; equipos heterogéneos},
      3/{selección}/{criterios explícitos; buscar casos que contradigan; registrar quién se negó y por qué},
      4/{memoria}/{anclar en fechas y eventos; contrastar con documentos y otros relatos},
      5/{decir vs. hacer}/{no tratar el relato como hecho: analizarlo como discurso situado; triangular con observación},
      6/{interpretación impuesta}/{preguntas interpretativas; control con los participantes (\emph{member checking})}} {
      \node[r] (r\i) at (0,{-0.95*\i}) {\textbf{\a}};
      \node[s] (s\i) at (6.2,{-0.95*\i}) {\b};
      \draw[arr] (r\i) -- (s\i);
    }
  \end{tikzpicture}}
  \end{center}
  \begin{center}
  \tiny \textcite[cap.~23]{vieytes2004metodologia}; Maxwell, en \textcite[cap.~3]{bickmanrog1998handbook}; \textcite{rinaldoguhin2022interviews}; \textcite{smallcalarco2022qualitative}.
  \end{center}
\end{frame}

% ============================ SECCIÓN 4 ============================
\section{Variantes de entrevista}

\begin{frame}[t]{Una tipología en dos criterios}
  \begin{center}
  \adjustbox{max width=0.95\textwidth, max height=0.6\textheight}{%
  \begin{tikzpicture}[cell/.style={bx, text width=2.8cm, minimum height=1.0cm, font=\tiny}]
    \node[font=\scriptsize\bfseries] at (0,1.0) {estructurada};
    \node[font=\scriptsize\bfseries] at (3.2,1.0) {semiestructurada};
    \node[font=\scriptsize\bfseries] at (6.4,1.0) {no estructurada};
    \node[font=\scriptsize, text=utdtGray] at (3.2,1.55) {grado de espontaneidad de la interacción $\longrightarrow$};
    \foreach \i/\r in {0/{cara a cara}, 1/{telefónica}, 2/{virtual}} {
      \node[font=\scriptsize\bfseries, anchor=east] at (-1.6,{-1.2*\i}) {\r};
      \node[cell, fill=utdtGray!10] at (0,{-1.2*\i}) {cuestionario de encuesta};
      \node[cell, fill=utdtYellow] at (3.2,{-1.2*\i}) {en profundidad};
      \node[cell, fill=utdtYellow] at (6.4,{-1.2*\i}) {en profundidad};
    }
    \draw[decorate, decoration={brace, amplitude=6pt, mirror}, thick] (1.7,-3.0) -- (7.9,-3.0) node[midway, below=8pt, font=\scriptsize]{lo que hoy llamamos entrevista \textbf{cualitativa}};
  \end{tikzpicture}}
  \end{center}
  \begin{center}
  \tiny Fideli y Marradi (1996), ampliado por Piovani en \textcite[tabla~12.1, p.~217]{marradi2007metodologia}. Las grupales (\emph{focus group}) se agregan como otra variante: el dato es la \textbf{interacción} \parencite{cyr2016focus}.
  \end{center}
\end{frame}

\begin{frame}[t]{Variantes y sus principales atributos}
  \begin{center}
  \adjustbox{max width=0.98\textwidth, max height=0.7\textheight}{%
  \begin{tikzpicture}[h/.style={font=\tiny\bfseries, text width=3.2cm, align=center}, c/.style={bxg, font=\tiny, text width=3.2cm, minimum height=0.62cm}, k/.style={bxk, font=\tiny\bfseries, text width=2.3cm, minimum height=0.62cm}]
    \foreach \x/\t in {2.85/{rasgo central}, 6.3/{sirve para}, 9.75/{cuidado}} { \node[h] at (\x,0.6) {\t}; }
    \foreach \i/\v/\a/\b/\c in {
      0/{semiestructurada}/{guía de temas; orden y redacción flexibles}/{comparar entre entrevistados}/{no convertir la guía en cuestionario},
      1/{no estructurada}/{consigna amplia; el entrevistado lleva el hilo}/{explorar marcos de sentido}/{difícil de comparar},
      2/{focalizada}/{todos vivieron un mismo estímulo o situación}/{experiencia subjetiva de un hecho}/{definir bien el estímulo},
      3/{historia de vida}/{toda la trayectoria, con documentos personales}/{articular biografía e historia}/{memoria; extensión},
      4/{historia oral}/{eventos o períodos específicos del pasado}/{dar voz a grupos marginados}/{olvidos y distorsiones},
      5/{grupal / focus group}/{4 a 12 personas y un moderador}/{normas y discursos colectivos}/{voces dominantes; consenso forzado},
      6/{a élites y expertos}/{personas con poder o saber técnico}/{decisiones y procesos internos}/{asimetría; discurso oficial},
      7/{en línea}/{video sincrónica o correo asincrónico}/{dispersión; costo; flexibilidad}/{menos \emph{rapport}; brecha digital}} {
      \node[k] at (0,{-0.72*\i}) {\v};
      \node[c] at (2.85,{-0.72*\i}) {\a};
      \node[c] at (6.3,{-0.72*\i}) {\b};
      \node[c] at (9.75,{-0.72*\i}) {\c};
    }
  \end{tikzpicture}}
  \end{center}
  \begin{center}
  \tiny \textcite{mertonkendall1946focused}; \textcite[secs.~19.5--19.6 y cap.~20]{brymanetal2021socialresearch}; Bornat, en \textcite[cap.~20]{alasuutarietal2008handbook}; \textcite[caps.~14, 21 y 23]{vieytes2004metodologia}.
  \end{center}
\end{frame}

\begin{frame}[t]{¿Presencial o en línea? Lo que muestra la evidencia}
  \begin{center}
  \adjustbox{max width=0.95\textwidth, max height=0.6\textheight}{%
  \begin{tikzpicture}
    \node[bx, text width=3.2cm, minimum height=1.4cm] (p) at (0,0) {\textbf{presencial}\\más empatía y \emph{rapport} percibidos};
    \node[bxk, text width=3.2cm, minimum height=1.4cm] (z) at (7.0,0) {\textbf{videollamada}\\más flexible, menor costo, llega lejos};
    \draw[biarr, draw=utdtGray] (p) -- node[lbl, above=2pt]{misma cantidad de palabras}  node[lbl, below=2pt]{(calidad similar)} (z);
    \node[bxg, font=\tiny, text width=3.2cm] (p2) at (0,-2.0) {preferible con temas que generan \textbf{vergüenza} y con personas mayores};
    \node[bxg, font=\tiny, text width=3.2cm] (z2) at (7.0,-2.0) {los más \textbf{jóvenes} reportaron más \emph{rapport} por video};
    \draw[garr] (p2) -- (p); \draw[garr] (z2) -- (z);
    \node[bxr, font=\tiny, text width=2.8cm] (r) at (3.5,-2.0) {por correo: respuestas más pensadas, menos espontáneas};
  \end{tikzpicture}}
  \end{center}
  \begin{center}
  \tiny \textcite{anthonyetal2025inperson} (73 participantes, diseño cuasiexperimental); \textcite{oliffeetal2021zoom}; Curasi (2001), en \textcite[sec.~19.5]{brymanetal2021socialresearch}.
  \end{center}
\end{frame}

\begin{frame}[t]{Entrevistas en investigaciones publicadas}
  \begin{center}
  \adjustbox{max width=0.98\textwidth, max height=0.68\textheight}{%
  \begin{tikzpicture}[e/.style={bx, text width=3.6cm, minimum height=1.9cm, font=\tiny}]
    \node[e, fill=utdtYellow!40] at (0,0) {\textbf{Small y Cook (2023)}, \emph{SMR}\\[2pt] cinco desafíos al preguntar ``por qué'' y estrategias para estudiar la acción motivada\\[2pt]\emph{semiestructurada}};
    \node[e] at (4.0,0) {\textbf{Rinaldo y Guhin (2022)}, \emph{SMR}\\[2pt] por qué funcionan las entrevistas: acceso a sentido, emociones y relatos situados\\[2pt]\emph{defensa metodológica}};
    \node[e, fill=utdtYellow!40] at (8.0,0) {\textbf{Anthony et al. (2025)}, \emph{IJQM}\\[2pt] compara presencial y videollamada en \emph{rapport}, empatía y extensión\\[2pt]\emph{en línea}};
    \node[e] at (0,-2.3) {\textbf{Khan y Jerolmack (2013)}, \emph{Sociological Quarterly}\\[2pt] estudiantes de élite: lo que dicen y lo que hacen\\[2pt]\emph{entrevista vs. observación}};
    \node[e, fill=utdtYellow!40] at (4.0,-2.3) {\textbf{Pérez (2018)}, \emph{Mobilization}\\[2pt] trayectorias de militantes desocupados reconstruidas en relatos de vida\\[2pt]\emph{historia de vida}};
    \node[e] at (8.0,-2.3) {\textbf{Lamont y Swidler (2014)}, \emph{Qualitative Sociology}\\[2pt] contra el ``tribalismo metodológico''\\[2pt]\emph{qué pueden y qué no pueden las entrevistas}};
  \end{tikzpicture}}
  \end{center}
  \begin{center}
  \tiny \textcite{smallcook2023why}; \textcite{rinaldoguhin2022interviews}; \textcite{anthonyetal2025inperson}; \textcite{khanjerolmack2013saying}; \textcite{perez2018piquetero}; \textcite{lamontswidler2014pluralism}. Guía práctica: \textcite{lareau2021listening}.
  \end{center}
\end{frame}

% ============================ EJERCICIO 1 ============================
\section{Ejercicio: preguntas de entrevista}

\begin{frame}[t]{Ejercicio: corregir una guía}
  \begin{columns}[T]
    \begin{column}{0.56\textwidth}
      \scriptsize
      Tema: \textbf{cómo organizan su vida los estudiantes que trabajan}. Un borrador de guía:
      \vspace{0.1cm}
      \begin{center}
      \adjustbox{max width=\textwidth}{%
      \begin{tikzpicture}[q/.style={doc, text width=6.4cm}]
        \node[q] at (0,0) {\textbf{a.} ¿Trabaja?};
        \node[q] at (0,-0.6) {\textbf{b.} ¿No cree que trabajar perjudica sus notas?};
        \node[q] at (0,-1.2) {\textbf{c.} ¿Cómo concilia la conciliación trabajo--estudio y su capital cultural?};
        \node[q] at (0,-1.8) {\textbf{d.} ¿Le gusta su trabajo y le paga bien su jefe?};
        \node[q] at (0,-2.4) {\textbf{e.} ¿Por qué decidió trabajar?};
        \node[q] at (0,-3.0) {\textbf{f.} ¿Ha pensado en abandonar la carrera? (primera pregunta)};
      \end{tikzpicture}}
      \end{center}
    \end{column}
    \begin{column}{0.42\textwidth}
      \begin{center}
      \adjustbox{max width=\textwidth, max height=0.66\textheight}{%
      \begin{tikzpicture}[s/.style={bx, text width=4.2cm, minimum height=0.7cm}]
        \node[s, fill=utdtYellow] (a) at (0,0) {\textbf{1.} ¿Qué problema tiene cada pregunta?};
        \node[s] (b) at (0,-1.25) {\textbf{2.} Reescríbanlas como preguntas abiertas y neutras};
        \node[s] (c) at (0,-2.5) {\textbf{3.} ¿Qué tipo de Kvale es cada versión nueva?};
        \node[s] (d) at (0,-3.75) {\textbf{4.} Ordénenlas en \textbf{embudo} y agreguen una pregunta de sondeo};
        \node[s, fill=utdtGray!10] (e) at (0,-5.0) {\textbf{5.} ¿Qué parte del tema conviene \textbf{observar} en vez de preguntar?};
        \foreach \x/\y in {a/b, b/c, c/d, d/e} { \draw[arr] (\x) -- (\y); }
      \end{tikzpicture}}
      \end{center}
    \end{column}
  \end{columns}
\end{frame}

\begin{frame}[t]{Puesta en común: una lectura posible}
  \begin{center}
  \adjustbox{max width=0.97\textwidth, max height=0.7\textheight}{%
  \begin{tikzpicture}[q/.style={bxg, font=\tiny, text width=3.2cm, minimum height=0.8cm}, a/.style={bx, font=\tiny, text width=7.6cm, minimum height=0.8cm, align=left}]
    \foreach \i/\x/\y in {
      0/{a. cerrada (sí/no)}/{``Cuénteme cómo es una semana normal para usted, entre el trabajo y la universidad.'' (introductoria)},
      1/{b. dirigida}/{``¿Cómo se relacionan, en su experiencia, el trabajo y sus estudios?'' (indirecta del juicio)},
      2/{c. jerga y redundancia}/{``¿Qué hace cuando un turno coincide con una evaluación? Recuerde la última vez.'' (especificación)},
      3/{d. doble pregunta}/{dividir: ``¿Qué le gusta y qué no de su trabajo?'' y, aparte, condiciones laborales},
      4/{e. invita a justificar}/{``¿Cómo llegó a su primer trabajo? ¿Qué pasaba en su vida en ese momento?'' (relato, no motivo abstracto)},
      5/{f. directa y sensible al inicio}/{moverla al final; antes, sondeo: ``Mencionó que el semestre pasado fue difícil, ¿qué ocurrió?''}} {
      \node[q] at (0,{-0.95*\i}) {\x};
      \node[a] at (5.6,{-0.95*\i}) {\y};
    }
    \node[bxk, font=\tiny, text width=11.0cm] at (2.8,-6.0) {\textbf{5.} los horarios, los traslados y el uso del tiempo de estudio se pueden \textbf{observar} o reconstruir con un diario solicitado; la entrevista conviene para el \textbf{sentido} que le dan};
  \end{tikzpicture}}
  \end{center}
\end{frame}

% ============================ SECCIÓN 6 ============================
\section{Información secundaria cualitativa}

\begin{frame}[t]{¿Qué es un documento para la investigación?}
  \begin{center}
  \adjustbox{max width=0.97\textwidth, max height=0.6\textheight}{%
  \begin{tikzpicture}
    \node[bxk, text width=2.6cm, minimum height=1.5cm, font=\scriptsize\bfseries] (d) at (0,0) {documento};
    \node[bx, text width=2.8cm] (a) at (-5.0,1.3) {se puede \textbf{leer} (también imágenes)};
    \node[bx, text width=2.8cm] (b) at (5.0,1.3) {\textbf{no fue producido} para la investigación};
    \node[bx, text width=2.8cm] (c) at (-5.0,-1.3) {fue \textbf{preservado} y está disponible};
    \node[bx, text width=2.8cm] (e) at (5.0,-1.3) {es \textbf{relevante} para la pregunta};
    \draw[arr] (a) -- (d); \draw[arr] (b) -- (d); \draw[arr] (c) -- (d); \draw[arr] (e) -- (d);
    \node[bxg, font=\scriptsize, text width=4.2cm] (nr) at (0,-2.5) {ventaja: \textbf{no reactivo}; nuestra recolección no altera el dato};
    \draw[darr] (d) -- (nr);
    \node[bxr, font=\tiny, text width=3.4cm] (nf) at (0,2.4) {pero no es más fácil: buscar, reunir e interpretar cuesta tiempo};
  \end{tikzpicture}}
  \end{center}
  \begin{center}
  \tiny \textcite[secs.~22.1--22.2]{brymanetal2021socialresearch}. Para Piovani, es \textbf{observación indirecta}: conocemos el objeto a través de observaciones hechas por otros \parencite[p.~193]{marradi2007metodologia}.
  \end{center}
\end{frame}

\begin{frame}[t]{Un mapa de las fuentes secundarias cualitativas}
  \begin{center}
  \adjustbox{max width=0.98\textwidth, max height=0.66\textheight}{%
  \begin{tikzpicture}[n/.style={bx, text width=2.3cm, minimum height=0.9cm}, l/.style={bxg, font=\tiny, text width=2.3cm, minimum height=1.1cm}]
    \node[bx, fill=utdtDark, text=white, text width=2.6cm] (r) at (6.0,0) {\textbf{fuentes secundarias cualitativas}};
    \node[n, fill=utdtYellow] (p) at (0,-1.5) {\textbf{personales}};
    \node[n] (e) at (3.0,-1.5) {\textbf{oficiales del Estado}};
    \node[n] (o) at (6.0,-1.5) {\textbf{oficiales privadas}};
    \node[n] (m) at (9.0,-1.5) {\textbf{medios masivos}};
    \node[n] (v) at (12.0,-1.5) {\textbf{digitales}};
    \node[l] at (0,-3.0) {cartas, diarios, autobiografías, fotos familiares};
    \node[l] at (3.0,-3.0) {leyes, informes, actas, expedientes, archivos públicos};
    \node[l] at (6.0,-3.0) {memorias, actas de directorio, manuales, correos internos};
    \node[l] at (9.0,-3.0) {prensa, televisión, cine, publicidad};
    \node[l] at (12.0,-3.0) {sitios, blogs, foros, publicaciones en redes};
    \foreach \x in {p,e,o,m,v} { \draw[arr] (r) -- (\x); }
    \node[bxk, font=\tiny, text width=4.4cm] (q) at (6.0,-4.6) {\textbf{datos cualitativos de otros}: entrevistas y notas de campo archivadas para su reutilización};
  \end{tikzpicture}}
  \end{center}
  \begin{center}
  \tiny Tipología de \textcite{scott1990matter}, en \textcite[cap.~22]{brymanetal2021socialresearch}; ``documentos de vida'' de \textcite{plummer1983documents}; reutilización de datos archivados: Heaton, en \textcite[cap.~30]{alasuutarietal2008handbook}.
  \end{center}
\end{frame}

\begin{frame}[t]{Procedimiento: de la pregunta al corpus}
  \begin{center}
  \adjustbox{max width=0.98\textwidth, max height=0.6\textheight}{%
  \begin{tikzpicture}[s/.style={bx, text width=1.9cm, minimum height=1.7cm, font=\tiny}]
    \node[s, fill=utdtYellow] (a) at (0,0) {\textbf{1. pregunta}\\delimita qué buscar};
    \node[s] (b) at (2.35,0) {\textbf{2. identificar}\\archivos, bibliotecas, bases, repositorios};
    \node[s] (c) at (4.7,0) {\textbf{3. acceder}\\reglas del archivo; solicitudes de información pública};
    \node[s] (d) at (7.05,0) {\textbf{4. delimitar}\\el corpus: período, lugar, tipo; muestreo};
    \node[s] (e) at (9.4,0) {\textbf{5. registrar}\\fichas y \textbf{bitácora} de búsqueda};
    \node[s] (f) at (11.75,0) {\textbf{6. evaluar}\\críticamente cada fuente};
    \node[s] (g) at (14.1,0) {\textbf{7. analizar}\\contenido, discurso, semiótica};
    \foreach \x/\y in {a/b, b/c, c/d, d/e, e/f, f/g} { \draw[arr] (\x) -- (\y); }
    \node[bxg, font=\scriptsize, text width=9.0cm] (q) at (7.05,-2.1) {``las fuentes jamás hablan por sí mismas'': revisar bases de datos esperando que algo nos sorprenda solo produce angustia};
    \draw[darr] (q) -- (a.south);
  \end{tikzpicture}}
  \end{center}
  \begin{center}
  \tiny \textcite[cap.~6, pp.~266--273]{vieytes2004metodologia}; identificación, acceso, selección y análisis de archivos en \textcite[concepto clave 22.1]{brymanetal2021socialresearch}.
  \end{center}
\end{frame}

\begin{frame}[t]{El documento como recurso y como tema}
  \begin{center}
  \adjustbox{max width=0.95\textwidth, max height=0.56\textheight}{%
  \begin{tikzpicture}
    \node[bx, text width=3.4cm, minimum height=1.6cm] (r) at (0,0) {\textbf{como recurso}\\contenedor de información sobre otra cosa (qué pasó, quién, cuándo)};
    \node[bxk, text width=3.4cm, minimum height=1.6cm] (t) at (7.4,0) {\textbf{como tema}\\cómo se produjo, circula y qué \textbf{hace} en la vida social};
    \draw[biarr, draw=utdtGray] (r) -- (t);
    \node[bxg, font=\tiny, text width=3.2cm] at (0,-2.0) {usar actas de una cooperativa para fechar decisiones};
    \node[bxg, font=\tiny, text width=3.2cm] at (7.4,-2.0) {estudiar cómo esas actas construyen la legitimidad de la directiva};
  \end{tikzpicture}}
  \end{center}
  \begin{center}
  \tiny Prior, en \textcite[cap.~28]{alasuutarietal2008handbook}. Ejemplo de documentos como tema: los informes oficiales que presentan un barrio como ``problema'' para justificar una gran obra pública (Thompson et al., 2013, en \cite[sec.~22.4]{brymanetal2021socialresearch}).
  \end{center}
\end{frame}

\begin{frame}[t]{Cuatro criterios para evaluar un documento}
  \begin{center}
  \adjustbox{max width=0.97\textwidth, max height=0.64\textheight}{%
  \begin{tikzpicture}[c/.style={bx, text width=3.4cm, minimum height=1.5cm}, xn/.style={bxg, font=\tiny, text width=3.2cm}]
    \node[bxk, circle, minimum size=1.8cm, text width=1.6cm, font=\scriptsize\bfseries] (s) at (0,0) {calidad de la fuente};
    \node[c] (a) at (-5.0,1.6) {\textbf{autenticidad}\\¿es genuino y de origen indudable?};
    \node[c] (b) at (5.0,1.6) {\textbf{credibilidad}\\¿está libre de error y distorsión?};
    \node[c] (c) at (-5.0,-1.6) {\textbf{representatividad}\\¿es típico de su clase? ¿se sabe cuán atípico es?};
    \node[c] (d) at (5.0,-1.6) {\textbf{significado}\\¿es claro y comprensible en su contexto?};
    \foreach \n in {a,b,c,d} { \draw[arr] (\n) -- (s); }
    \node[xn] at (-5.0,3.0) {autobiografías con escritores fantasma};
    \node[xn] at (5.0,3.0) {cartas de famosos que saben que serán leídas};
    \node[xn] at (-5.0,-3.0) {escribían los alfabetizados: élites y varones};
    \node[xn] at (5.0,-3.0) {lo que el autor y el destinatario no necesitaban decir};
  \end{tikzpicture}}
  \end{center}
  \begin{center}
  \tiny \textcite[p.~6]{scott1990matter}, en \textcite[sec.~22.2]{brymanetal2021socialresearch}. Los criterios pesan distinto según el tipo: un documento estatal suele ser auténtico, pero no necesariamente creíble ni representativo.
  \end{center}
\end{frame}

\begin{frame}[t]{Las fotografías tampoco hablan solas}
  \begin{columns}[T]
    \begin{column}{0.5\textwidth}
      \begin{center}
      \adjustbox{max width=\textwidth, max height=0.6\textheight}{%
      \begin{tikzpicture}[s/.style={bx, text width=4.0cm, minimum height=0.8cm, font=\tiny}]
        \node[font=\scriptsize\bfseries] at (0,0.7) {tres tipos de foto familiar};
        \node[s, fill=utdtYellow!50] at (0,0) {\textbf{idealización}: pose formal (matrimonio, graduación)};
        \node[s] at (0,-1.0) {\textbf{retrato natural}: la acción mientras ocurre, con algo de puesta en escena};
        \node[s] at (0,-2.0) {\textbf{desmitificación}: situaciones atípicas o incómodas};
        \node[font=\scriptsize\bfseries] at (0,-3.0) {tres usos en investigación};
        \node[s, fill=utdtGray!10] at (0,-3.7) {\textbf{ilustración} $\cdot$ \textbf{dato} $\cdot$ \textbf{disparador} en entrevistas};
      \end{tikzpicture}}
      \end{center}
    \end{column}
    \begin{column}{0.48\textwidth}
      \scriptsize
      \begin{itemize}
        \item ``Toffs and toughs'' (Londres, 1937) se usa como imagen de la desigualdad de clases: los dos jóvenes de sombrero de copa esperaban un auto para una fiesta, y los otros tres no eran ``matones'', sino chicos del barrio que ganaban propinas.
        \item Importa tanto lo \textbf{fotografiado} como lo que \textbf{no} se fotografió.
      \end{itemize}
      \vspace{0.1cm}
      \tiny \textcite{scott1990matter}; \textcite[sec.~22.3]{brymanetal2021socialresearch}.
    \end{column}
  \end{columns}
\end{frame}

% ============================ SECCIÓN 7 ============================
\section{Solo secundaria o como complemento}

\begin{frame}[t]{¿Cuándo trabajar con datos disponibles?}
  \begin{center}
  \adjustbox{max width=0.98\textwidth, max height=0.62\textheight}{%
  \begin{tikzpicture}[c/.style={bx, text width=3.0cm, minimum height=1.1cm}]
    \node[bxk, circle, minimum size=2.0cm, text width=1.8cm, font=\scriptsize\bfseries] (o) at (0,0) {datos disponibles};
    \node[c] (c1) at (-5.0,1.7) {los hechos \textbf{ya ocurrieron} (ex post facto)};
    \node[c] (c2) at (0,2.4) {fenómenos \textbf{poco frecuentes}, difíciles de presenciar};
    \node[c] (c3) at (5.0,1.7) {no es \textbf{ético} producir el dato de primera mano};
    \node[c] (c4) at (-5.0,-1.7) {falta de \textbf{tiempo} o excesiva \textbf{distancia}};
    \node[c] (c5) at (0,-2.4) {la investigación primaria es demasiado \textbf{costosa}};
    \node[c] (c6) at (5.0,-1.7) {existen datos \textbf{válidos y confiables} sobre lo que interesa};
    \foreach \c in {c1,c2,c3,c4,c5,c6} { \draw[arr] (\c) -- (o); }
  \end{tikzpicture}}
  \end{center}
  \begin{center}
  \tiny \textcite[cap.~6, sec.~6.5.1, p.~275]{vieytes2004metodologia}. En la investigación histórica, a menudo es la \textbf{única} vía.
  \end{center}
\end{frame}

\begin{frame}[t]{Cuando solo hay información secundaria}
  \begin{center}
  \adjustbox{max width=0.98\textwidth, max height=0.66\textheight}{%
  \begin{tikzpicture}[m/.style={bx, text width=2.2cm, minimum height=1.6cm, font=\tiny}]
    \node[font=\scriptsize\bfseries] at (5.2,1.3) {cinco modos de reutilizar datos cualitativos de otros};
    \node[m, fill=utdtYellow] at (0,0) {\textbf{supra-análisis}\\preguntas nuevas que trascienden el estudio original};
    \node[m] at (2.6,0) {\textbf{suplementario}\\un tema que el estudio original trató poco};
    \node[m, fill=utdtYellow] at (5.2,0) {\textbf{re-análisis}\\verificar hallazgos con otro marco};
    \node[m] at (7.8,0) {\textbf{amplificado}\\combinar dos o más bases};
    \node[m, fill=utdtYellow] at (10.4,0) {\textbf{surtido}\\combinar con datos nuevos o documentos};
    \node[bxr, font=\tiny, text width=4.2cm] (pb) at (1.6,-2.2) {problema: \textbf{no haber estado allí}; se pierde el conocimiento tácito del campo};
    \node[bxk, font=\tiny, text width=5.4cm] (rs) at (8.2,-2.2) {respuesta: \textbf{metadocumentación} (guías, notas, memos), contacto con el equipo original, revisar el \textbf{consentimiento} para reutilizar};
    \draw[arr] (pb) -- (rs);
  \end{tikzpicture}}
  \end{center}
  \begin{center}
  \tiny Heaton, en \textcite[cap.~30]{alasuutarietal2008handbook}. \textcite{kernmustasilta2023beyond} muestran que el análisis secundario cualitativo no ahorra tiempo, pero amplía la comparación de casos y reduce la fatiga de comunidades muy investigadas.
  \end{center}
\end{frame}

\begin{frame}[t]{Recomendaciones cuando solo se cuenta con documentos}
  \begin{center}
  \adjustbox{max width=0.97\textwidth, max height=0.66\textheight}{%
  \begin{tikzpicture}[s/.style={bx, text width=9.6cm, minimum height=0.68cm, font=\tiny, align=left}, k/.style={bxk, font=\tiny\bfseries, text width=0.5cm, minimum height=0.68cm}]
    \foreach \i/\t in {
      1/{\textbf{reconstruir el contexto de producción}: quién escribió, para quién, con qué propósito y bajo qué restricciones},
      2/{\textbf{buscar fuentes contrapuestas}: el documento oficial junto a la prensa opositora, las cartas privadas o los testimonios},
      3/{\textbf{controlar explicaciones alternativas}: los datos que permitirían descartarlas pueden no haberse registrado nunca},
      4/{\textbf{documentar la búsqueda}: bitácora de archivos, palabras clave, lo encontrado y lo que no se pudo consultar},
      5/{\textbf{declarar el alcance}: conclusiones sobre lo que los documentos permiten afirmar, no sobre todo lo ocurrido},
      6/{\textbf{reconocer el ``no haber estado allí''} y, si es posible, sumar entrevistas a quienes produjeron o conservaron las fuentes}} {
      \node[k] at (0,{-0.85*\i}) {\i};
      \node[s] at (5.3,{-0.85*\i}) {\t};
    }
  \end{tikzpicture}}
  \end{center}
  \begin{center}
  \tiny \textcite[cap.~6, secs.~6.5.2 y 6.6]{vieytes2004metodologia}; \textcite[cap.~22]{brymanetal2021socialresearch}; \textcite{mayrlwilson2020historical}, sobre cómo los sociólogos históricos organizan su trabajo con archivos.
  \end{center}
\end{frame}

\begin{frame}[t]{Cuando la información secundaria es un complemento}
  \begin{center}
  \adjustbox{max width=0.97\textwidth, max height=0.62\textheight}{%
  \begin{tikzpicture}
    \node[bxk, text width=2.8cm, minimum height=1.4cm] (p) at (0,0) {\textbf{datos primarios}\\observación y entrevistas};
    \node[bx, text width=2.8cm, minimum height=1.4cm] (s) at (10.0,0) {\textbf{datos secundarios}\\documentos, archivos, datos ajenos};
    \draw[arr] (s.north west) to[bend right=18] node[lbl, above]{\textbf{antes}: contexto, selección de casos, guía de entrevista} (p.north east);
    \draw[arr] (s.west) -- node[lbl]{\textbf{durante}: disparadores (fotos, cartas) en la entrevista} (p.east);
    \draw[arr] (s.south west) to[bend left=18] node[lbl, below]{\textbf{después}: contrastar lo dicho y lo observado} (p.south east);
    \node[bxg, font=\scriptsize, text width=8.4cm] at (5.0,-2.7) {el documento no ``verifica'' mecánicamente el relato: cada fuente tiene su propio sesgo y la discrepancia también es un dato};
  \end{tikzpicture}}
  \end{center}
  \begin{center}
  \tiny Triangulación con ``documentos de vida'' \parencite{plummer1983documents}, en \textcite[p.~201]{marradi2007metodologia}; las historias de vida combinan relato con diarios, fotos y cartas \parencite[sec.~19.6]{brymanetal2021socialresearch}.
  \end{center}
\end{frame}

% ============================ SECCIÓN 8 ============================
\section{Sesgos en la producción de datos secundarios}

\begin{frame}[t]{Antes de llegar a nuestras manos: los silencios del archivo}
  \begin{center}
  \adjustbox{max width=0.98\textwidth, max height=0.6\textheight}{%
  \begin{tikzpicture}[s/.style={bx, text width=2.6cm, minimum height=1.5cm}, r/.style={bxr, font=\tiny, text width=2.6cm}]
    \node[s] (a) at (0,0) {\textbf{1. creación}\\quién pudo producir un registro};
    \node[s] (b) at (3.4,0) {\textbf{2. conservación}\\qué se guardó en un archivo};
    \node[s] (c) at (6.8,0) {\textbf{3. recuperación}\\qué se busca y se encuentra};
    \node[s, fill=utdtYellow] (d) at (10.2,0) {\textbf{4. interpretación}\\qué relato construimos};
    \foreach \x/\y in {a/b, b/c, c/d} { \draw[arr] (\x) -- (\y); }
    \node[r] at (0,-1.8) {alfabetizados, élites, Estado; los subalternos dejan menos huellas};
    \node[r] at (3.4,-1.8) {\textbf{supervivencia selectiva}: daño, descarte, destrucción};
    \node[r] at (6.8,-1.8) {catálogos, buscadores y palabras clave; acceso restringido};
    \node[r] at (10.2,-1.8) {elegir solo los documentos que confirman (\emph{cherry picking})};
  \end{tikzpicture}}
  \end{center}
  \begin{center}
  \tiny Momentos de producción de silencios en \textcite{trouillot1995silencing}; ``selección diferencial por la historia'' en \textcite{kingkeohaneverba1994designing}; supervivencia selectiva de cartas y fotos en \textcite[sec.~22.3]{brymanetal2021socialresearch}.
  \end{center}
\end{frame}

\begin{frame}[t]{Preguntas críticas antes de usar un dato ajeno}
  \begin{center}
  \adjustbox{max width=0.98\textwidth, max height=0.68\textheight}{%
  \begin{tikzpicture}[q/.style={bx, text width=4.6cm, minimum height=0.9cm, font=\tiny}]
    \node[q, fill=utdtYellow] at (0,0) {¿Tengo acceso al \textbf{proceso de producción} del dato original?};
    \node[q] at (0,-1.1) {¿Es \textbf{creíble} la fuente que lo produjo? ¿Quién la financió?};
    \node[q, fill=utdtYellow] at (0,-2.2) {¿Quién fue el \textbf{destinatario original}? Se escribe distinto para un juez, un jefe o un amigo};
    \node[q] at (0,-3.3) {¿Cuál es el \textbf{contexto temporal}? Evitar lecturas anacrónicas};
    \node[q] at (5.4,0) {¿Coincide la \textbf{unidad de análisis} con la mía?};
    \node[q, fill=utdtYellow] at (5.4,-1.1) {¿Qué \textbf{universo} cubre el corpus y qué queda fuera?};
    \node[q] at (5.4,-2.2) {¿Me interesan los \textbf{aspectos} del concepto que el documento registra?};
    \node[q, fill=utdtYellow] at (5.4,-3.3) {¿Puedo \textbf{contrastarlo} con otras fuentes o repetir parte de la recolección?};
  \end{tikzpicture}}
  \end{center}
  \begin{center}
  \tiny Adaptado de la evaluación crítica de datos secundarios de \textcite[cap.~6, secs.~6.6.1--6.6.13, pp.~278--287]{vieytes2004metodologia}, pensada para datos numéricos y útil también para documentos cualitativos.
  \end{center}
\end{frame}

\begin{frame}[t]{Procedimientos para evitar sesgos}
  \begin{center}
  \adjustbox{max width=0.97\textwidth, max height=0.7\textheight}{%
  \begin{tikzpicture}[r/.style={bxr, text width=3.2cm, minimum height=0.8cm, font=\tiny}, s/.style={bxk, text width=6.8cm, minimum height=0.8cm, font=\tiny}]
    \node[font=\scriptsize\bfseries] at (0,0.75) {sesgo};
    \node[font=\scriptsize\bfseries] at (6.2,0.75) {procedimiento};
    \foreach \i/\a/\b in {
      0/{autoría desigual}/{buscar deliberadamente voces ausentes: archivos sindicales, de organizaciones, testimonios orales},
      1/{supervivencia selectiva}/{reconstruir la historia del archivo: qué se descartó, cuándo y por qué},
      2/{selección del investigador}/{definir el corpus \textbf{antes} de leerlo; criterios explícitos; incluir casos negativos},
      3/{destinatario y propósito}/{leer cada documento junto a su género (informe, carta, publicidad) y su audiencia},
      4/{buscadores y digitalización}/{registrar consultas; probar sinónimos y términos de época; ir también al papel},
      5/{anacronismo}/{glosario de términos de época; leer documentos contemporáneos que den contexto},
      6/{un solo tipo de fuente}/{triangular tipos de documento y, si es posible, sumar entrevistas}} {
      \node[r] (r\i) at (0,{-0.92*\i}) {\textbf{\a}};
      \node[s] (s\i) at (6.2,{-0.92*\i}) {\b};
      \draw[arr] (r\i) -- (s\i);
    }
  \end{tikzpicture}}
  \end{center}
  \begin{center}
  \tiny \textcite[sec.~22.3 y aprender de la experiencia 22.2]{brymanetal2021socialresearch}; \textcite[cap.~6]{vieytes2004metodologia}; \textcite{trouillot1995silencing}; \textcite{kernmustasilta2023beyond}.
  \end{center}
\end{frame}

\begin{frame}[t]{Documentos en investigaciones publicadas}
  \begin{center}
  \adjustbox{max width=0.98\textwidth, max height=0.68\textheight}{%
  \begin{tikzpicture}[e/.style={bx, text width=3.6cm, minimum height=1.9cm, font=\tiny}]
    \node[e, fill=utdtYellow!40] at (0,0) {\textbf{Nelson (2021)}, \emph{AJS}\\[2pt] un siglo de documentos de movimientos de mujeres en dos ciudades, con análisis computacional e interpretativo\\[2pt]\emph{archivos organizacionales}};
    \node[e] at (4.0,0) {\textbf{Mayrl y Wilson (2020)}, \emph{AJS}\\[2pt] cómo los sociólogos históricos seleccionan y combinan fuentes de archivo\\[2pt]\emph{reflexión sobre el método}};
    \node[e, fill=utdtYellow!40] at (8.0,0) {\textbf{Kern y Mustasilta (2023)}, \emph{CPS}\\[2pt] análisis secundario de entrevistas y grupos focales ajenos; recomendaciones de transparencia\\[2pt]\emph{reutilización}};
    \node[e] at (0,-2.3) {\textbf{Edelmann et al. (2020)}, \emph{Annual Review of Sociology}\\[2pt] huellas digitales y textos masivos como nuevas fuentes\\[2pt]\emph{documentos digitales}};
    \node[e, fill=utdtYellow!40] at (4.0,-2.3) {\textbf{Nelson (2020)}, \emph{SMR}\\[2pt] teoría fundamentada computacional: lectura humana y algoritmos sobre grandes corpus\\[2pt]\emph{puente con la semana 12}};
    \node[e] at (8.0,-2.3) {\textbf{Ejemplos en Bryman}\\[2pt] cartas de un soldado de 1862 y el blog de otro en Irak; informes oficiales sobre un barrio ``a regenerar''\\[2pt]\emph{personales y estatales}};
  \end{tikzpicture}}
  \end{center}
  \begin{center}
  \tiny \textcite{nelson2021cycles}; \textcite{mayrlwilson2020historical}; \textcite{kernmustasilta2023beyond}; \textcite{edelmannetal2020computational}; \textcite{nelson2020computational}; Shapiro y Humphreys (2013) y Thompson et al. (2013), en \textcite[cap.~22]{brymanetal2021socialresearch}.
  \end{center}
\end{frame}

% ============================ EJERCICIO 2 ============================
\section{Ejercicio: evaluar un documento}

\begin{frame}[t]{Ejercicio: una carta encontrada en un archivo}
  \begin{columns}[T]
    \begin{column}{0.54\textwidth}
      \scriptsize
      Pregunta: \textbf{¿cómo vivieron los trabajadores textiles el cierre de su fábrica?} En el archivo de la empresa aparece:
      \vspace{0.1cm}
      \begin{center}
      \begin{tikzpicture}
        \node[doc, text width=6.0cm] {\textit{12 de marzo de 1979.} A todo el personal: La Gerencia lamenta comunicar que, por la situación del mercado, la planta cesará sus actividades el 30 de abril. Agradecemos la comprensión y el espíritu de colaboración que el personal ha demostrado siempre. Todos los trabajadores recibirán lo que corresponde según la ley. \textit{Firma ilegible. Copia mecanografiada, sin membrete.}};
      \end{tikzpicture}
      \end{center}
      \tiny Documento ficticio, elaborado para el ejercicio.
    \end{column}
    \begin{column}{0.44\textwidth}
      \begin{center}
      \adjustbox{max width=\textwidth, max height=0.66\textheight}{%
      \begin{tikzpicture}[s/.style={bx, text width=4.4cm, minimum height=0.7cm}]
        \node[s, fill=utdtYellow] (a) at (0,0) {\textbf{1.} Evalúen \textbf{autenticidad}, \textbf{credibilidad}, \textbf{representatividad} y \textbf{significado}};
        \node[s] (b) at (0,-1.35) {\textbf{2.} ¿Quién escribe, para quién y con qué propósito?};
        \node[s] (c) at (0,-2.6) {\textbf{3.} ¿Recurso o tema? ¿Qué puede y qué no puede responder?};
        \node[s] (d) at (0,-3.85) {\textbf{4.} ¿Qué \textbf{voces faltan} y dónde podrían buscarlas?};
        \node[s, fill=utdtGray!10] (e) at (0,-5.1) {\textbf{5.} ¿Con qué otra técnica lo \textbf{complementarían}?};
        \foreach \x/\y in {a/b, b/c, c/d, d/e} { \draw[arr] (\x) -- (\y); }
      \end{tikzpicture}}
      \end{center}
    \end{column}
  \end{columns}
\end{frame}

\begin{frame}[t]{Puesta en común: una lectura posible}
  \begin{center}
  \adjustbox{max width=0.97\textwidth, max height=0.7\textheight}{%
  \begin{tikzpicture}[q/.style={bxg, font=\scriptsize, text width=2.6cm, minimum height=0.85cm}, a/.style={bx, font=\tiny, text width=8.6cm, minimum height=0.85cm, align=left}]
    \foreach \i/\x/\y in {
      0/{1. Scott}/{\textbf{autenticidad} dudosa (copia, sin membrete, firma ilegible): buscar el original o registros de envío; \textbf{credibilidad} baja sobre causas (``situación del mercado''); \textbf{representatividad}: ¿se conservaron otras comunicaciones?; \textbf{significado}: ``lo que corresponde'' depende de la ley de la época},
      1/{2. producción}/{la gerencia, para el personal y también para la posteridad (y quizás para la justicia): busca mostrar el cierre como inevitable y legítimo},
      2/{3. recurso o tema}/{como \textbf{recurso}: fecha y forma del anuncio; como \textbf{tema}: el lenguaje de la ``colaboración''. No dice cómo lo \textbf{vivieron} los trabajadores},
      3/{4. voces ausentes}/{archivo sindical, prensa local, expedientes judiciales laborales, fotos y cartas familiares},
      4/{5. complemento}/{entrevistas de \textbf{historia oral} a ex trabajadores, usando la carta como \textbf{disparador}; contrastar memorias con fechas documentadas}} {
      \node[q] at (0,{-1.15*\i}) {\x};
      \node[a] at (6.0,{-1.15*\i}) {\y};
    }
  \end{tikzpicture}}
  \end{center}
\end{frame}

% ============================ CIERRE ======================================
\section{Cierre}

\begin{frame}[t]{Mapa de las dos clases}
  \begin{center}
  \adjustbox{max width=0.98\textwidth, max height=0.68\textheight}{%
  \begin{tikzpicture}[t/.style={bxk, text width=2.6cm, minimum height=1.0cm}, a/.style={bx, font=\tiny, text width=2.6cm, minimum height=0.85cm}]
    \node[t] (o) at (0,0) {\textbf{observación}\\hacer};
    \node[t] (e) at (4.6,0) {\textbf{entrevista}\\decir};
    \node[t] (d) at (9.2,0) {\textbf{documentos}\\huellas};
    \foreach \n/\x/\c/\p/\s in {
      o/0/{prácticas, lo tácito, lo oculto}/{acceso, rol, embudo, notas}/{reactividad, selectividad},
      e/4.6/{sentidos, pasado, temas sensibles}/{selección, guía, Kvale, tácticas}/{preguntas dirigidas, deseabilidad},
      d/9.2/{hechos pasados, sin reacción}/{corpus, bitácora, Scott}/{supervivencia, autoría, destinatario}} {
      \node[a] (c\n) at (\x,-1.4) {\textbf{cuándo}: \c};
      \node[a] (p\n) at (\x,-2.45) {\textbf{cómo}: \p};
      \node[a, draw=BrickRed] (s\n) at (\x,-3.5) {\textbf{sesgo}: \s};
      \draw[garr] (\n) -- (c\n);
    }
    \draw[biarr, dashed, draw=utdtGray] (o) -- node[lbl]{triangulación} (e);
    \draw[biarr, dashed, draw=utdtGray] (e) -- node[lbl]{triangulación} (d);
  \end{tikzpicture}}
  \end{center}
\end{frame}

\begin{frame}{Checklist para su proyecto}
  \footnotesize
  \begin{itemize}
    \item[$\square$] ¿Justifiqué por qué \textbf{entrevistar} (y no solo observar) para responder mi pregunta?
    \item[$\square$] ¿Tengo criterios explícitos para \textbf{seleccionar} entrevistados y una regla para decidir cuántos?
    \item[$\square$] ¿Mi guía va en \textbf{embudo}, con preguntas abiertas, neutras y sin jerga? ¿La probé?
    \item[$\square$] ¿Qué \textbf{documentos} existen sobre mi caso, quién los produjo y para quién?
    \item[$\square$] ¿Evalué mis fuentes con los criterios de \textbf{Scott} y registré mi búsqueda en una \textbf{bitácora}?
    \item[$\square$] ¿Qué \textbf{voces faltan} en mis datos y cómo intentaré incorporarlas?
  \end{itemize}
\end{frame}

\begin{frame}[t]{Próxima semana}
  \begin{center}
  \adjustbox{max width=0.9\textwidth, max height=0.5\textheight}{%
  \begin{tikzpicture}[s/.style={bx, text width=3.8cm, minimum height=1.6cm}]
    \node[s, fill=utdtGray!10] (a) at (0,0) {\textbf{Semana 11}\\producir datos: observar, entrevistar, reunir documentos};
    \node[s, fill=utdtYellow] (b) at (5,0) {\textbf{Semana 12}\\analizar datos cualitativos: codificación y teoría fundamentada};
    \draw[arr] (a) -- (b);
  \end{tikzpicture}}
  \end{center}
  \begin{center}
  \footnotesize\itshape Lectura: \textcite{timmermanstavory2012theory}.
  \end{center}
\end{frame}

\begin{frame}[standout]
  ¿Preguntas?

  \vspace{0.4cm}
  {\small \href{mailto:\CorreoProfesor}{\CorreoProfesor}}
\end{frame}

\begin{frame}[allowframebreaks]{Referencias}
  \printbibliography[heading=none]
\end{frame}

\end{document}
